% Faithful formal-publication transcription of complete Section3.2 proof.
\section*{3.2 Proof of Theorem3}
We conclude this overview section with a proof of Theorem3 using Multileaf Lemma. We start with a lemma that shows how the noisy leaves in the domain of $g^n$ can be ``simulated'' by low-degree conical juntas in the domain of $\mathrm{GapOr}^n$. For simplicity, we state the lemma assuming a hard side $b=0$; an analogous lemma holds for $b=1$ by replacing GapOr with $\neg\mathrm{GapOr}$.
\begin{authorclaim}{Simulation Lemma}
Let $\ell$ be a $(\delta,\epsilon,\mu,0)$-noisy leaf over the variables of $g^n$. There exists a conical junta $p_\ell:(\{0,1\}^{\log n})^n\to\mathbb R_{\ge0}$ of degree at most $n\cdot[\delta\log n+\log(1/\epsilon)]$ such that
\[\forall x\in\operatorname{dom}(\mathrm{GapOr}_{\log n}^n):\quad p_\ell(x)=\ell(\mu^{\mathrm{GapOr}_{\log n}^n(x)}).\]
\end{authorclaim}
\begin{proof}
We start by defining three conical juntas in the domain of $\mathrm{GapOr}_m$ for $m:=\log n$. Let $S_k^m$ be the distribution over multisets obtained by picking $k$ random elements from $[m]$ with replacement.
\begin{align*}
q_1(x)&:=\tfrac2m\sum_{i\in[m]}x_i&&\text{of degree }1,\\
q_2(x)&:=\prod_{i\in[m]}\overline x_i&&\text{of degree }m=\log n,\\
q_3(x)&:=\mathbb E_{S\sim S_k^m}\prod_{i\in S}\overline x_i&&\text{of degree }k:=\log(1/\epsilon).
\end{align*}
Note the following output values:
\begin{align*}
\forall x\in(\mathrm{GapOr}_m)^{-1}(0):&\quad q_1(x)=0,\quad q_2(x)=1,\quad q_3(x)=1,\\
\forall x\in(\mathrm{GapOr}_m)^{-1}(1):&\quad q_1(x)=1,\quad q_2(x)=0,\quad q_3(x)=2^{-k}=\epsilon.
\end{align*}
Let $y=(y^1,\ldots,y^n)$ be the input variables of $g^n$. We write $\ell(y)=\prod_i\ell_i(y^i)$ so that $\ell(\mu^{\mathrm{GapOr}_m^n(x)})=\prod_i\ell_i(\mu^{\mathrm{GapOr}_m(x^i)})$. We simulate each factor in this product separately. For $i\in[n]$ consider the function $p_i:\{0,1\}^m\to\mathbb R_{\ge0}$ defined by
\[p_i(x):=\ell_i(\mu^{\mathrm{GapOr}_m(x)}).\]
First note that $p_i$ can always be written as a conical combination of $q_1$ and $q_2$ in degree $\log n$. Moreover, if $\ell_i$ is $(\epsilon,\mu,0)$-noisy, meaning $\ell_i(\mu^1)\ge\epsilon\cdot\ell_i(\mu^0)$, then we can do better and write $p_i$ as a conical combination of $q_1$ and $q_3$ in degree $\log(1/\epsilon)$. We now define $p_\ell:=\prod_ip_i$. The claimed bound on the degree of $p_\ell$ follows because at most $\delta$ fraction of the $\ell_i$ are non-noisy.
\end{proof}
We are now ready to prove Theorem3 using Multileaf Lemma and Simulation Lemma.
\begin{proof}[Proof of Theorem3]
Suppose for contradiction that $T$ is a randomised decision tree for $f\circ g^n$ having error $1/3$ and depth $\gamma n\overline R_{1/n}(g)$ where $\gamma=o(1)$ as $n\to\infty$. Our goal is to construct an $o(n\log n)$-degree $o(1)$-approximate conical junta for $f\circ\mathrm{GapOr}_{\log n}^n$ (or $f\circ\neg\mathrm{GapOr}_{\log n}^n$).
We make two simplifying assumptions wlog.
\begin{enumerate}
\item The randomised tree $T$ has error $o(1)$. To ensure this, we may reduce $T$'s error by running it $1/\sqrt\gamma=\omega(1)$ times. This will yield an $o(1)$-error tree of depth $\sqrt\gamma n\overline R_{1/n}(g)=o(n\overline R_{1/n}(g))$.
\item There is some $\epsilon:=1/n^{o(1)}$ such that $T$ has depth $o(n\overline R_\epsilon(g))$. To ensure this, we may apply Claim7 to see that $\gamma n\overline R_{1/n}(f)\le\sqrt\gamma n\overline R_\epsilon(f)\le o(n\overline R_\epsilon(f))$ where $\epsilon:=1/n^{4\sqrt\gamma}$.
\end{enumerate}
We invoke Multileaf Lemma with the above $\epsilon\le o(1)$ and $\delta:=\max\{\gamma^c,\epsilon^c\}\le o(1)$ for small enough constant $c>0$. We get a hard distribution $\mu$ and a hard side $b$, say $b=0$ (case $b=1$ is similar, but using $\neg\mathrm{GapOr}$), such that the following holds: For every $t$ in the support of $T$ if we run $t$ on a random input $x\sim\mu^y$, where $y\in\{0,1\}^n$, then the leaf reached $\ell_x^t$ will be $(\delta,\epsilon,\mu,0)$-noisy with probability $1-o(1)$. This allows us to effectively ignore non-noisy leaves: denoting by $N(t)\subseteq A(t)$ the set of accepting leaves that are $(\delta,\epsilon,\mu,0)$-noisy, we have
\begin{equation}\tag{2}
\begin{aligned}
\forall y\in\{0,1\}^n:\quad\mathbb E_{x\sim\mu^y}[t(x)]
&=\mathbb E_{x\sim\mu^y}\left[\sum_{\ell\in A(t)}\ell(x)\right]&&\text{(Using(1))}\\
&\in\mathbb E_{x\sim\mu^y}\left[\sum_{\ell\in N(t)}\ell(x)\right]\pm o(1).
\end{aligned}
\end{equation}
We now define the approximating conical junta by
\[p(x):=\mathbb E_{t\sim T}\left[\sum_{\ell\in N(t)}p_\ell(x)\right],\]
where the $p_\ell$ are given by Simulation Lemma. Hence $p$ has degree at most
\[n\cdot[\delta\log n+\log(1/\epsilon)]=n\cdot[o(1)\log n+\log n^{o(1)}]=o(n\log n).\]
We finish the proof of Theorem3 by verifying that $p$ indeed $o(1)$-approximates $f\circ\mathrm{GapOr}_{\log n}^n$.
\begin{align*}
\forall x:\quad p(x)&=\mathbb E_{t\sim T}\left[\sum_{\ell\in N(t)}p_\ell(x)\right]&&\bigl(y:=\mathrm{GapOr}_{\log n}^n(x)\bigr)\\
&=\mathbb E_{t\sim T}\left[\sum_{\ell\in N(t)}\ell(\mu^y)\right]\\
&=\mathbb E_{t\sim T}\left[\sum_{\ell\in N(t)}\mathbb E_{x'\sim\mu^y}[\ell(x')]\right]\\
&=\mathbb E_{t\sim T}\left[\mathbb E_{x'\sim\mu^y}\left[\sum_{\ell\in N(t)}\ell(x')\right]\right]\\
&\in\mathbb E_{t\sim T}\left[\mathbb E_{x'\sim\mu^y}[t(x')]\right]\pm o(1)&&\text{(Using(2))}\\
&=\mathbb E_{x'\sim\mu^y}\left[\mathbb E_{t\sim T}[t(x')]\right]\pm o(1)\\
&\in\mathbb E_{x'\sim\mu^y}[(f\circ g^n)(x')]\pm o(1)&&\text{($T$ has error $o(1)$)}\\
&=f(y)\pm o(1)=(f\circ\mathrm{GapOr}_{\log n}^n)(x)\pm o(1).
\end{align*}
\end{proof}
